% Source: MQ36_v1.html
\documentclass[11pt]{article}
\usepackage[utf8]{inputenc}
\usepackage[T1]{fontenc}
\usepackage[margin=1in]{geometry}
\usepackage{amsmath}
\usepackage{amssymb}
\usepackage{siunitx}
\usepackage{enumitem}
\usepackage{parskip}

\title{PS160 -- MQ36\_v1.html}
\date{}

\begin{document}
\maketitle

\section*{Questions}

\begin{enumerate}[leftmargin=*]
  \item A single slit experiment involves monochromatic, coherent light incident on an aperture of width 0.26 mm. If the screen is 1.8 m away, and the distance from the central maximum to the \emph{m}= 5 minimum is 7.3 mm, determine the wavelength of the light in nanometers. \hfill \textit{(1 pts, Numeric)}

  \item The light emitted from an excited atom has two principal wavelengths: 589.3 nm and 588.52 nm. If you wish to resolve these two emission lines in first order using a diffraction grating that is 1.667 cm in length, what minimum number of slits per centimeter must the grating have? \hfill \textit{(1 pts, Numeric)}

  \item X-rays of wavelength 0.22 nm incident at a grazing angle of 22.7$^{\circ}$ on an unknown crystal produces a first-order interference fringe. Calculate the interatomic spacing between the Bragg planes in nanometers. \hfill \textit{(1 pts, Numeric)}

  \item Your pupil has a diameter of 5 mm. Using 551 nm light, what is the resolving power of your eye? Express your answer in milli-degrees. \hfill \textit{(1 pts, Numeric)}

\end{enumerate}

\end{document}